\documentclass[12pt]{article}
\usepackage[utf8]{inputenc}
\usepackage[T1]{fontenc}
\usepackage[margin=1in]{geometry}
\usepackage{setspace}
\usepackage{xcolor}
\usepackage[colorlinks=true, urlcolor=blue!60!black, linkcolor=blue!60!black]{hyperref}
\setlength{\parindent}{2em}
\setlength{\parskip}{1.25em}
\renewcommand{\baselinestretch}{1.0}
\begin{document}
\begin{center}
{\LARGE \textit{Jackson v. Metropolitan Edison Co.}, 419 U.S. 345 (1974)\par}
\vspace{0.8em}
{\large Case Brief --- Four-Page Edition\par}
\end{center}
\vspace{1em}

\section{Memory Jogger}

A privately owned, heavily regulated electric utility with at least a partial monopoly that cut off a customer's power without notice or a hearing was not a state actor: the Court required a close nexus between the State and the specific challenged act and confined the public-function doctrine to powers traditionally exclusively reserved to the State (\href{https://supreme.justia.com/cases/federal/us/419/345/}{419 U.S.\ 345}).

\section{Detailed Case Facts}

Metropolitan Edison, a privately owned Pennsylvania corporation, held a certificate of public convenience to supply electricity to an area including York and was extensively regulated by the state Public Utility Commission (PUC). Rule~15 of its PUC-filed tariff reserved the right to discontinue service on reasonable notice for nonpayment.

Catherine Jackson's account was terminated in 1970 for delinquency, and service resumed under an account for James Dodson, another occupant. After Dodson moved out in 1971, service continued without payment. On October 7, 1971, an employee told Jackson the meter had been tampered with; she denied knowledge and asked that the account be put in the name of ``Robert Jackson,'' her 12-year-old son, and, as Justice Brennan noted, offered \$30 toward her old bill. On October 11, without further notice, Metropolitan disconnected her service.

\section{Procedural History}

Jackson sued under 42 U.S.C.\ \S\ 1983, claiming that a state-law entitlement to reasonably continuous service could not be ended without notice and a hearing. The Middle District of Pennsylvania dismissed for lack of state action, and the Third Circuit affirmed. Citing a circuit split, the Supreme Court granted certiorari, heard argument on October 15, 1974, and affirmed on December 23, 1974.

\section{Judicial Votes}

Affirmed, \textbf{6--3}.
\begin{itemize}
\item \textbf{Majority:} Justice Rehnquist, joined by Chief Justice Burger and Justices Stewart, White, Blackmun, and Powell.
\item \textbf{Dissenting:} Justices Douglas, Brennan, and Marshall, each separately.
\end{itemize}

\section{Holding}

Pennsylvania was not sufficiently connected with the termination to make it state action. A plaintiff must show a ``sufficiently close nexus'' between the State and the specific challenged action of the regulated entity. The public-function doctrine reaches only powers ``traditionally exclusively reserved to the State,'' which does not include furnishing utility service. The Court did not decide whether continued service was a protected property interest.

\section{Analysis of Opinions}

\textbf{Majority (Rehnquist).} Starting from the \textit{Civil Rights Cases}, Rehnquist placed the case on the \textit{Moose Lodge} side of the line rather than \textit{Burton}'s, rejecting four theories. \emph{Regulation:} extensive regulation alone is not enough; otherwise ``a lone Philadelphia cab driver'' would be the State. \emph{Monopoly:} utilities are natural monopolies, and in any event monopoly status had no relationship to the termination. \emph{Public function:} unlike elections, company towns, municipal parks, or eminent domain, supplying electricity has never been the State's exclusive prerogative, and ``affected with a public interest'' would sweep in doctors, lawyers, and a grocer selling milk. \emph{Tariff approval:} Rule~15 took effect automatically without PUC scrutiny; unlike \textit{Pollak}, where the commission investigated and endorsed the practice, the State had not put its weight behind the shutoff. Nor was there a \textit{Burton}-style symbiotic relationship: the State was not ``a partner or even a joint venturer.''

\textbf{Douglas (dissenting).} The majority's method was ``fundamentally sequential, rather than cumulative''; in the aggregate, licensing, monopoly, regulation, and tariff empowerment made the shutoff state action, and commissions had become ``prisoners of the utilities.''

\textbf{Brennan (dissenting).} Jackson was not a customer of record, so no controversy existed; he would have vacated and remanded for dismissal.

\textbf{Marshall (dissenting).} Pennsylvania foreclosed competition and delegated eminent domain, making service the functional equivalent of the State's; requiring the PUC to \emph{order} a practice gutted \textit{Pollak}; and pluralism values are irrelevant for the only electric company in town, so the holding would seemingly shield racial refusals of service.

\textbf{Later developments.} \textit{Flagg Bros.}, \textit{Blum v.\ Yaretsky}, \textit{Rendell-Baker}, and \textit{Lugar} extended \textit{Jackson}'s limits, and \textit{Manhattan Community Access Corp.\ v.\ Halleck} (2019) relied on it to hold that running public-access channels is not an exclusive public function.

\section{Examples: Future Applications}

\begin{enumerate}
\item \textbf{Same-Side (No State Action): Nursing-Home Transfer.} A licensed, Medicaid-funded nursing home's physician-driven transfer is private, as \textit{Blum v.\ Yaretsky} held.
\item \textbf{Same-Side (No State Action): Internet Service.} A franchised internet provider's shutoff under unreviewed filed terms lacks a nexus, and internet service is no exclusive state function.
\item \textbf{Opposite-Side (State Action): Commission-Ordered Shutoff.} A commission order directing utilities to disconnect meters without a hearing supplies state compulsion.
\item \textbf{Opposite-Side (Likely State Action): Eminent Domain.} A utility condemning an easement under delegated power presents the case \textit{Jackson} said ``would be quite a different one.''
\item \textbf{Fence-Sitter: Charter Schools.} A fully state-funded charter school's expulsion sits between \textit{Rendell-Baker} and \textit{Brentwood}'s entwinement theory; the Fourth Circuit in \textit{Peltier} found state action.
\end{enumerate}

\section{Critique}

\textbf{Core critiques.} Critics fault the seriatim dismissal of factors that Douglas would have weighed together; the legal-realist problem of treating utility self-help and commission silence as neutral private action despite state-law backing; and the weak pluralism rationale for shielding an investor-owned monopoly.

\textbf{The formalist cliff.} In \textit{Memphis Light, Gas \& Water Division v.\ Craft} (1978), a municipally owned utility's shutoff on similar facts required pre-termination notice and a chance to contest billing errors. Identical shutoffs thus trigger due process or none at all depending solely on ownership.

\section{Key Quotations}

\begin{quote}
``But the inquiry must be whether there is a sufficiently close nexus between the State and the challenged action of the regulated entity so that the action of the latter may be fairly treated as that of the State itself.''
\end{quote}
\noindent --- Justice Rehnquist, for the Court, 419 U.S.\ at 351.

\begin{quote}
``We have, of course, found state action present in the exercise by a private entity of powers traditionally exclusively reserved to the State. \dots\ If we were dealing with the exercise by Metropolitan of some power delegated to it by the State which is traditionally associated with sovereignty, such as eminent domain, our case would be quite a different one.''
\end{quote}
\noindent --- Justice Rehnquist, for the Court, 419 U.S.\ at 352--353.

\begin{quote}
``Approval by a state utility commission of such a request from a regulated utility, where the commission has not put its own weight on the side of the proposed practice by ordering it, does not transmute a practice initiated by the utility and approved by the commission into `state action.'\,''
\end{quote}
\noindent --- Justice Rehnquist, for the Court, 419 U.S.\ at 357.

\end{document}
